%% pc-telemetre-ultrasons — mesure d'une distance par écho ultrasonore.
%% Haut : le montage (capteur É/R, obstacle, cote d).
%% Bas  : le chronogramme (salve émise, écho reçu, durée aller-retour Delta t).
\documentclass[tikz,border=4pt]{standalone}
\usepackage{liaisons}
\begin{document}
\begin{tikzpicture}[x=1cm,y=1cm,
  onde/.style={-{Stealth[length=6pt]}, line width=1.2pt},
  cote/.style={{Stealth[length=4pt]}-{Stealth[length=4pt]}, solideD, line width=0.9pt},
  axe/.style={-{Stealth[length=5pt]}, line width=0.6pt},
  rappel/.style={line width=0.4pt, dash pattern=on 1.5pt off 1.5pt, black!55},
  txt/.style={font=\scriptsize}]

%% ======================= partie haute : le montage ========================
%% sol
\draw[line width=0.8pt] (-0.2,0) -- (6.1,0);
\foreach \i in {-0.1,0.2,...,6.0} { \draw[line width=0.4pt, black!60] (\i,0) -- (\i-0.14,-0.14); }
%% capteur
\fill[black!12] (0,0.3) rectangle (1.1,1.4);
\draw[line width=0.7pt] (0,0.3) rectangle (1.1,1.4);
\node[txt, anchor=south, align=center] at (0.55,1.45) {capteur};
\foreach \yy/\lb in {1.1/{É}, 0.6/{R}} {
  \draw[line width=0.6pt, fill=white] (0.92,\yy) circle (0.16);
  \node[font=\tiny] at (0.92,\yy) {\lb}; }
%% obstacle
\hachures{(5.3,0) rectangle (5.8,1.65)}
\node[txt, anchor=south] at (5.55,1.7) {obstacle};
%% salve émise et écho
\draw[onde, solideA] (1.2,1.1) -- (5.2,1.1) node[txt, midway, above, inner sep=2pt, text=solideA] {salve émise};
\draw[onde, solideB] (5.2,0.62) -- (1.2,0.62) node[txt, midway, below, inner sep=2pt, text=solideB] {écho};
%% cote d au niveau du sol
\draw[cote] (0.55,0.14) -- (5.3,0.14)
   node[midway, fill=white, inner sep=1pt, font=\small, text=solideD] {$d$};
\draw[rappel] (0.55,0.05) -- (0.55,0.3);

%% ======================= partie basse : le chronogramme ===================
\begin{scope}[shift={(0,-2.35)}]
  \draw[axe] (0.5,0) -- (6.1,0) node[txt, anchor=north east, inner sep=3pt] {$t$ (s)};
  \draw[axe] (0.5,-0.1) -- (0.5,1.6) node[txt, anchor=south west, inner sep=1pt] {tension (V)};
  %% pic d'émission (haut) et pic de l'écho (plus bas, plus loin)
  \draw[solideA, line width=1pt, line join=round] (0.6,0) -- (0.8,1.1) -- (1.0,0);
  \draw[solideB, line width=1pt, line join=round] (3.7,0) -- (3.9,0.55) -- (4.1,0);
  \node[txt, anchor=north, text=solideA, align=center] at (0.8,-0.05) {émission};
  \node[txt, anchor=north, text=solideB, align=center] at (3.9,-0.05) {réception\\de l'écho};
  %% durée aller-retour
  \draw[rappel] (0.8,1.13) -- (0.8,1.47);
  \draw[rappel] (3.9,0.58) -- (3.9,1.47);
  \draw[cote] (0.8,1.4) -- (3.9,1.4)
     node[midway, fill=white, inner sep=1pt, font=\small, text=solideD] {$\Delta t$};
\end{scope}

%% ======================= relation encadrée ================================
\node[draw, line width=0.7pt, rounded corners=2pt, inner sep=4pt, align=center,
      anchor=north] at (3.1,-3.2)
  {$d = \dfrac{v_{\text{US}} \times \Delta t}{2}$\\[2pt]
   {\scriptsize pendant $\Delta t$, l'onde parcourt $2d$ : \textbf{aller-retour}}};
\end{tikzpicture}
\end{document}
